\section{鞅：随信息演化的条件平均}\label{sec:06-02}

\subsection{滤过、鞅与条件平均的时间一致性}\label{subsec:6-2-1}
\index{滤过}\index{鞅}\index{次鞅}\index{可预测过程}

设 $X\in L^1$ 是一天结束时才能完全看见的赔付，而
$\mathcal F_0\subset\mathcal F_1\subset\cdots$ 是逐日收到的信息。时刻 $n$ 的自然估计是
\[
  M_n=\E[X\mid\mathcal F_n].
\]
如果到了时刻 $m<n$，先计算第 $n$ 日的估计再退回第 $m$ 日，塔式性质给出
\begin{equation}\label{eq:6-2-time-consistency-opening}
  \E[M_n\mid\mathcal F_m]
  =\E[\E[X\mid\mathcal F_n]\mid\mathcal F_m]
  =\E[X\mid\mathcal F_m]=M_m.
\end{equation}
这条时间一致性比“平均不变”强得多：它要求每一层已有信息上的条件平均都不变。鞅正是把
\eqref{eq:6-2-time-consistency-opening} 从一个终端变量的逐步估计中抽出来的结构。

\begin{definition}[滤过与适应性]\label{def:6-2-1-1}
设 $(\Omega,\mathcal F,\Prob)$ 是概率空间。离散时间滤过是满足
\[
  \mathcal F_0\subset\mathcal F_1\subset\cdots\subset\mathcal F
\]
的子 $\sigma$-代数族；连续时间滤过 $(\mathcal F_t)_{t\ge0}$ 满足
$\mathcal F_s\subset\mathcal F_t$ 对所有 $s\le t$ 成立。

过程 $X=(X_t)$ 若对每个 $t$ 都有 $X_t$ 对 $\mathcal F_t$ 可测，称它对该滤过
\emph{适应}（adapted）。离散时间过程 $(H_n)_{n\ge1}$ 若 $H_n$ 对
$\mathcal F_{n-1}$ 可测，称它可预测（predictable）。换言之，$H_n$ 必须在第 $n$ 个增量
出现以前确定。
\end{definition}

\begin{definition}[鞅、次鞅与上鞅]\label{def:6-2-1-2}
设 $(X_n)$ 是实值、可积且适应的过程。若
\[
  \E[X_{n+1}\mid\mathcal F_n]=X_n\quad\text{a.s.},
\]
则称 $X$ 为鞅；若左边不小于 $X_n$，称为次鞅；若左边不大于 $X_n$，称为上鞅。
复值鞅仍用等式定义；次鞅与上鞅只用于实值过程。

连续时间过程 $X=(X_t)_{t\ge0}$ 称为鞅，是指每个 $X_t\in L^1$、过程适应，并且
\[
  \E[X_t\mid\mathcal F_s]=X_s\quad\text{a.s.}\qquad(0\le s\le t).
\]
连续时间次鞅和上鞅相应把等号改为 $\ge$ 和 $\le$。这里必须直接要求任意 $s\le t$；
连续时间没有可以迭代的“相邻一步”。
\end{definition}

\begin{definition}[自然滤过]\label{def:6-2-1-3}
过程 $X$ 的自然滤过为
\[
  \mathcal F_n^X=\sigma(X_0,\ldots,X_n)
  \quad\text{或}\quad
  \mathcal F_t^X=\sigma(X_s:0\le s\le t).
\]
它记录截至当前时刻可由过程本身观察到的全部事件。给过程加入额外信息会得到更大的滤过；
鞅性质依赖所选滤过，因而不能把滤过从记号中无声删去。
\end{definition}

\begin{remark}[连续时间滤过的通常条件]\label{remark:6-2-1-usual}
连续时间中常要求滤过完备且右连续：前者表示它包含终端 $\sigma$-代数中所有零集的子集，后者表示
\[
  \mathcal F_t=\bigcap_{s>t}\mathcal F_s.
\]
这两项合称\emph{通常条件}（usual conditions）。它们使几乎处处修改与停时极限具有合适的可测性。离散时间中不需要
右连续性；Brownian 自然滤过的完备右连续化见 \S\ref{subsec:6-4-1}。
\end{remark}

“martingale”原是赌博术语；Doob 把它从公平游戏的比喻提升为条件期望理论的中心对象。独立增量只是一种
产生鞅的机制，而“未来增量在现有信息下条件均值为零”才是保留下来的本质，因此似然比、调和函数和
后来的随机积分都能进入同一框架。

相邻一步的条件可由塔式性质迭代，从而恢复任意两个时刻之间的时间一致性。

\begin{proposition}[时间一致性]\label{proposition:6-2-1-1}
设 $(M_n)$ 可积且适应。它是鞅，当且仅当对所有整数 $0\le m\le n$，
\[
  \E[M_n\mid\mathcal F_m]=M_m\quad\text{a.s.}
\]
次鞅情形也有相应结论：一步条件等价于
$\E[X_n\mid\mathcal F_m]\ge X_m$ 对所有 $m\le n$ 成立；把次鞅结论用于 $-X$，便得到
上鞅版本。
\end{proposition}

\begin{proof}
任意时刻条件取 $n=m+1$ 便给出相邻一步条件。反过来，设一步条件成立。对 $m<n$，塔式性质与归纳假设给出
\[
 \E[M_n\mid\mathcal F_m]
 =\E\bigl[\E[M_n\mid\mathcal F_{n-1}]\mid\mathcal F_m\bigr]
 =\E[M_{n-1}\mid\mathcal F_m]=M_m.
\]
归纳从 $n=m+1$ 开始，故鞅结论成立。若 $X$ 是次鞅，则条件期望的单调性给
\[
 \E[X_n\mid\mathcal F_m]
 =\E\bigl[\E[X_n\mid\mathcal F_{n-1}]\mid\mathcal F_m\bigr]
 \ge \E[X_{n-1}\mid\mathcal F_m]\ge X_m,
\]
同样完成归纳。上鞅情形对 $-X$ 使用次鞅结论即可。
\end{proof}

\begin{example}[独立中心和]\label{ex:6-2-1-1}
设 $\xi_1,\xi_2,\ldots$ 独立、可积且 $\E\xi_k=0$，令
$S_0=0$、$S_n=\sum_{k=1}^n\xi_k$ 及
$\mathcal F_n=\sigma(\xi_1,\ldots,\xi_n)$。由于 $\xi_{n+1}$ 独立于
$\mathcal F_n$，
\[
 \E[S_{n+1}\mid\mathcal F_n]
 =S_n+\E[\xi_{n+1}\mid\mathcal F_n]
 =S_n+\E\xi_{n+1}=S_n.
\]
所以中心独立增量的部分和是鞅。若均值为 $\mu_k$，则
$S_n-\sum_{k\le n}\mu_k$ 才是鞅；适应性本身不会消去漂移。
\end{example}

\begin{example}[闭鞅]\label{ex:6-2-1-2}
给定 $X\in L^1$ 和递增信息 $(\mathcal F_n)$，令
$M_n=\E[X\mid\mathcal F_n]$。每个 $M_n$ 可积并对 $\mathcal F_n$ 可测，且
\[
  \E[M_{n+1}\mid\mathcal F_n]
  =\E[\E[X\mid\mathcal F_{n+1}]\mid\mathcal F_n]
  =\E[X\mid\mathcal F_n]=M_n.
\]
因此 $(M_n)$ 是鞅。它在每一时刻都是同一终端变量 $X$ 的条件平均，故称为由 $X$ 闭合的鞅。
\end{example}

\begin{example}[似然比鞅]\label{ex:6-2-1-3}
设 $P,Q$ 是同一可测空间上的概率，并且对每个 $n$，
$Q|_{\mathcal F_n}\ll P|_{\mathcal F_n}$。取有限阶段的 Radon--Nikodym 密度
\[
  L_n=\frac{\dd Q|_{\mathcal F_n}}{\dd P|_{\mathcal F_n}}.
\]
$L_n$ 非负、$\mathcal F_n$-可测且 $\E_P L_n=1$。对任意
$A\in\mathcal F_n$，
\[
 \int_A L_{n+1}\dd P=Q(A)=\int_A L_n\dd P.
\]
条件期望的刻画因而给出
$\E_P[L_{n+1}\mid\mathcal F_n]=L_n$。所以 $(L_n)$ 是非负 $P$-鞅。
这里的逐阶段绝对连续是假设，不是“观察同一份数据”自动带来的结论。
\end{example}

凸函数把平均守恒变成单向漂移。

\begin{theorem}[条件 Jensen 产生次鞅]\label{theorem:6-2-1-2}
若 $(M_n)$ 是实值鞅，$\varphi:\R\to\R$ 凸，并且每个
$\varphi(M_n)$ 可积，则 $(\varphi(M_n))$ 是次鞅。特别地，只要相应幂次可积，
$(|M_n|^p)$ 对 $p\ge1$ 是次鞅，$(M_n^2)$ 也是次鞅。
\end{theorem}

\begin{proof}
条件 Jensen 不等式（定理~\ref{theorem:5-2-3-1}）及鞅等式给出
\[
 \E[\varphi(M_{n+1})\mid\mathcal F_n]
 \ge \varphi\bigl(\E[M_{n+1}\mid\mathcal F_n]\bigr)
 =\varphi(M_n).
\]
$\varphi(M_n)$ 对 $\mathcal F_n$ 可测，故满足次鞅定义。函数
$x\mapsto|x|^p$（$p\ge1$）和 $x\mapsto x^2$ 均凸，代入即得两个特例。
\end{proof}

鞅也可由逐步产生的“新信息”来描述。

\begin{proposition}[鞅差分分解]\label{proposition:6-2-1-3}
令 $M=(M_n)_{n\ge0}$ 为鞅，并置 $D_n=M_n-M_{n-1}$（$n\ge1$）。则
$D_n$ 对 $\mathcal F_n$ 可测且
\[
  \E[D_n\mid\mathcal F_{n-1}]=0.
\]
反之，若可积变量 $D_n$ 对 $\mathcal F_n$ 可测并满足上述条件，则对任意
$M_0\in L^1(\mathcal F_0)$，
$M_n=M_0+\sum_{k=1}^nD_k$ 是鞅。若 $M_n\in L^2$，则不同差分在 $L^2$ 中正交：
$\E[D_j\overline{D_k}]=0$ 对 $j\ne k$ 成立。
\end{proposition}

\begin{proof}
适应性给出 $D_n$ 的 $\mathcal F_n$-可测性，而
\[
 \E[D_n\mid\mathcal F_{n-1}]
 =\E[M_n\mid\mathcal F_{n-1}]-M_{n-1}=0.
\]
反过来，有限和可积且适应，并且
\[
 \E[M_n\mid\mathcal F_{n-1}]
 =M_{n-1}+\E[D_n\mid\mathcal F_{n-1}]=M_{n-1}.
\]
最后设 $j<k$。此时 $D_j$ 对 $\mathcal F_j\subset\mathcal F_{k-1}$ 可测，条件期望取因子给出
\[
 \E[D_j\overline{D_k}]
 =\E\bigl[D_j\overline{\E[D_k\mid\mathcal F_{k-1}]}\bigr]=0.
\]
平方可积保证乘积可积，因此计算合法。
\end{proof}

“下注策略”只能使用下注前的信息；这正是可预测性的数学内容。

\begin{proposition}[可预测变换]\label{proposition:6-2-1-4}
设 $M$ 是鞅，$(H_k)_{k\ge1}$ 是有界可预测过程。定义
\[
 (H\mathbin{\boldsymbol\cdot}M)_0=0,
 \qquad
 (H\mathbin{\boldsymbol\cdot}M)_n
 =\sum_{k=1}^nH_k(M_k-M_{k-1}).
\]
则 $H\mathbin{\boldsymbol\cdot}M$ 是鞅。若 $X$ 是次鞅且 $H_k\ge0$，同一变换是次鞅。
\end{proposition}

\begin{proof}
令 $D_k=M_k-M_{k-1}$。有界性保证 $H_kD_k\in L^1$，而
$H_k$ 对 $\mathcal F_{k-1}$ 可测，故
\[
 \E[H_kD_k\mid\mathcal F_{k-1}]
 =H_k\E[D_k\mid\mathcal F_{k-1}]=0.
\]
命题~\ref{proposition:6-2-1-3} 的反向即说明变换为鞅。若 $X$ 是次鞅，
$\E[X_k-X_{k-1}\mid\mathcal F_{k-1}]\ge0$；再由 $H_k\ge0$ 得每个变换差分的条件
期望非负，所以所得过程是次鞅。
\end{proof}

可预测性是结论的必要结构条件。令 $M$ 为简单对称随机游走，若先看见
$D_k=M_k-M_{k-1}\in\{-1,1\}$ 再取 $H_k=\sgn(D_k)$，则
$H_kD_k=|D_k|=1$，累计收益恒为 $n$。这里 $H_k$ 对 $\mathcal F_k$ 可测，却不对
$\mathcal F_{k-1}$ 可测。这个系数使用了下注时尚不可得的结果，因而不是允许的可预测策略。

次鞅比鞅多出的系统漂移可以被唯一分离出来。

\begin{proposition}[离散 Doob 分解]\label{proposition:6-2-1-5}
设 $X=(X_n)_{n\ge0}$ 是可积次鞅。存在分解
\[
  X_n=M_n+A_n,
\]
其中 $M$ 是鞅，$A_0=0$，并且 $A$ 可积、可预测且逐点递增；该分解在不可分辨意义下唯一。具体地，
\begin{equation}\label{eq:6-2-doob-decomposition}
 A_n=\sum_{k=1}^n
 \E[X_k-X_{k-1}\mid\mathcal F_{k-1}],
 \qquad M_n=X_n-A_n.
\end{equation}
\end{proposition}

\begin{proof}
记
$\widetilde a_k=\E[X_k-X_{k-1}\mid\mathcal F_{k-1}]$。次鞅性给出
$\widetilde a_k\ge0$ a.s.。逐项置 $a_k=(\widetilde a_k)^+$；它仍是同一个条件期望的
$\mathcal F_{k-1}$-可测可积版本，并且对每个样本点都有 $a_k\ge0$。于是 $A_0=0$、
$A_n=\sum_{k\le n}a_k$ 可积、递增，并且对 $n\ge1$，$A_n$ 对
$\mathcal F_{n-1}$ 可测。再计算
\[
 \E[M_n-M_{n-1}\mid\mathcal F_{n-1}]
 =\E[X_n-X_{n-1}\mid\mathcal F_{n-1}]-a_n=0,
\]
故 $M$ 是鞅。

若还有 $X=M'+A'$，其中 $M'$ 为鞅且 $A'$ 满足所列条件，则
\[
 A'_n-A'_{n-1}
 =\E[X_n-X_{n-1}\mid\mathcal F_{n-1}],
\]
因为左边已对 $\mathcal F_{n-1}$ 可测，而 $M'$ 的增量条件均值为零。由 $A'_0=0$
逐项求和得到 $A'_n=A_n$ a.s.，继而 $M'_n=M_n$ a.s.。时间指标可数，故可把这些零集取并，
得到两组过程在同一个概率一事件上对所有 $n$ 同时相等。
\end{proof}

\begin{example}[平方随机游走的补偿]
设例~\ref{ex:6-2-1-1} 中的增量还满足
$\E\xi_k^2=\sigma^2<\infty$。由于 $S_{k-1}$ 对 $\mathcal F_{k-1}$ 可测，且
$\xi_k$ 独立于过去，
\begin{align*}
 \E[S_k^2-S_{k-1}^2\mid\mathcal F_{k-1}]
 &=\E[2S_{k-1}\xi_k+\xi_k^2\mid\mathcal F_{k-1}]\\
 &=\sigma^2.
\end{align*}
因此 $S_n^2$ 是次鞅，其 Doob 分解为
\[
 S_n^2=\underbrace{(S_n^2-n\sigma^2)}_{\text{鞅}}
       +\underbrace{n\sigma^2}_{\text{可预测补偿}}.
\]
同一计算将在 Poisson 过程中变成补偿计数，在 Brownian 理论中变成二次变差。
\end{example}

\begin{figure}[htbp]
\centering
\begin{tikzpicture}[node distance=12mm and 17mm]
  \node[book strong node] (f0) {$\mathcal F_0$};
  \node[book strong node,right=of f0] (f1) {$\mathcal F_1$};
  \node[book strong node,right=of f1] (f2) {$\mathcal F_2$};
  \node[book strong node,right=of f2] (fn) {$\cdots\subset\mathcal F_n$};
  \node[book warm node,above=15mm of f1] (x) {终端变量 $X$};
  \draw[book accent arrow] (f0)--node[below,book label] {信息增加}(f1);
  \draw[book accent arrow] (f1)--(f2);
  \draw[book accent arrow] (f2)--(fn);
  \draw[book dashed arrow] (x)--node[left,book label] {$\E[\,\cdot\mid\mathcal F_0]$}(f0);
  \draw[book dashed arrow] (x)--node[right,book label] {$M_1$}(f1);
  \draw[book dashed arrow] (x)--node[right,book label] {$M_2$}(f2);
  \draw[decorate,decoration={brace,mirror,amplitude=4pt},BookWarm]
    ($(f1.south west)+(0,-5mm)$)--($(f2.south east)+(0,-5mm)$)
    node[midway,below=5pt,book label] {差分只携带相邻层的新信息};
\end{tikzpicture}
\caption{终端变量向递增信息层作条件平均，形成闭鞅；塔式性质保证投影的时间一致性。}
\label{fig:6-2-1-filtration}
\end{figure}

适应只说明“目前能够读到 $X_n$”，并不说明下一步条件均值为零；可预测只说明策略在增量出现前
确定，也不保证其有界或可积。鞅定义中的可积性、适应性和条件平均三项缺一不可。连续时间还要另行
处理滤过的完备化与右连续化；图~\ref{fig:6-2-1-filtration} 只表达离散层级，不能替代这些条件。

\begin{exercise}
设 $X\in L^1$，$M_n=\E[X\mid\mathcal F_n]$。直接用条件期望的积分刻画而非塔式性质的结论，
验证 $M$ 的鞅等式。
\end{exercise}

\begin{exercise}
设 $M\in L^2$ 是实值离散鞅。证明
$\E M_n^2=\E M_0^2+\sum_{k=1}^n\E(M_k-M_{k-1})^2$，并指出每一步使用了
命题~\ref{proposition:6-2-1-3} 的哪一部分。
\end{exercise}

\begin{exercise}
在简单对称随机游走中取 $X_n=S_n+n$。验证 $X$ 适应且可积，但不是鞅；判断它是次鞅还是上鞅。
\end{exercise}

最大不等式将把终端时刻的矩控制转成整段路径的控制。为此必须允许“第一次越界”这样的随机时刻
进入证明。

\subsection{Doob 最大不等式、上穿与鞅收敛}\label{subsec:6-2-2}
\index{Doob 最大不等式}\index{上穿不等式}\index{鞅收敛定理}

逐个时刻估计 $\Prob(|M_k|\ge\lambda)$，再对 $k\le n$ 作并集估计，通常会多付出一个
随 $n$ 增长的因子。鞅结构允许在路径\emph{第一次}越界的时刻使用条件平均，把所有越界路径只计算
一次。另一方面，最大值有界仍不直接排除序列在两个水平之间永远振荡；上穿次数正是记录这种振荡的
离散计数器。

Doob 的最大不等式与收敛定理遵循同一个方法：先建立对所有有限时刻统一的估计，再让时间上界趋于
无穷。上穿论证则把一个可预测交易策略转化为路径不等式
\eqref{eq:6-2-upcross-pathwise}，从而将振荡次数与期望估计联系起来。

\begin{definition}[最大过程]\label{def:6-2-2-1}
对离散过程 $M$，定义
\[
  M_n^*=\max_{0\le k\le n}|M_k|.
\]
连续时间相应的表达式是 $\sup_{0\le s\le t}|M_s|$；在使用它以前还需保证路径可测性或右连续性，
本小节只讨论离散时间。
\end{definition}

\begin{definition}[上穿次数]\label{def:6-2-2-2}
给定实数列 $x_0,\ldots,x_n$ 与 $a<b$，区间 $[a,b]$ 的上穿次数
$U_n[a,b]$ 是满足
\[
  0\le s_1<t_1<s_2<t_2<\cdots<s_m<t_m\le n,
  \qquad x_{s_j}\le a,\quad x_{t_j}\ge b
\]
的最大整数 $m$。对随机过程 $X$，把 $x_k$ 换成 $X_k(\omega)$，便得到随机变量
$U_n[a,b]$。
\end{definition}

\begin{definition}[闭鞅与一致可积鞅]\label{def:6-2-2-3}
若存在 $Y\in L^1$ 使
\[
  M_n=\E[Y\mid\mathcal F_n]\quad\text{a.s.}
\]
对每个 $n$ 成立，则称 $M$ 为闭鞅。若随机变量族
$\{M_n:n\ge0\}$ 按定义~\ref{def:5-3-2-1} 一致可积，称 $M$ 为一致可积鞅。
\end{definition}

\begin{theorem}[Doob $L^1$ 最大不等式]\label{theorem:6-2-2-1}
设 $(X_k)_{0\le k\le n}$ 是非负次鞅。对每个 $\lambda>0$，
\begin{equation}\label{eq:6-2-doob-weak-one}
 \lambda\Prob\!\left(\max_{0\le k\le n}X_k\ge\lambda\right)
 \le \E\!\left[X_n\1_{\{\max_{k\le n}X_k\ge\lambda\}}\right]
 \le \E X_n.
\end{equation}
\end{theorem}

\begin{proof}
令
\[
 A_k=\{X_0<\lambda,\ldots,X_{k-1}<\lambda,\ X_k\ge\lambda\},
 \qquad 0\le k\le n,
\]
其中 $A_0=\{X_0\ge\lambda\}$。这些事件两两不交，$A_k\in\mathcal F_k$，且其并为
$A=\{\max_{k\le n}X_k\ge\lambda\}$。由命题~\ref{proposition:6-2-1-1} 的次鞅版本，
\begin{align*}
 \lambda\Prob(A_k)
 &\le \E[X_k\1_{A_k}]\\
 &\le \E\!\left[\E[X_n\mid\mathcal F_k]\1_{A_k}\right]
 =\E[X_n\1_{A_k}].
\end{align*}
对 $k$ 求和便得到第一个不等式。第二个不等式来自 $X_n\ge0$ 和
$\1_A\le1$。
\end{proof}

\begin{figure}[htbp]
\centering
\begin{tikzpicture}[node distance=7mm and 8mm]
  \node[book node] (a0) {$A_0$\\初始即越界};
  \node[book node,right=of a0] (a1) {$A_1$\\首次在 $1$ 越界};
  \node[right=of a1] (dots) {$\cdots$};
  \node[book node,right=of dots] (an) {$A_n$\\首次在 $n$ 越界};
  \node[book strong node,below=10mm of dots] (u)
    {$\displaystyle\bigsqcup_{k=0}^n A_k
      =\left\{\max_{j\le n}X_j\ge\lambda\right\}$};
  \draw[book arrow] (a0)--(u);
  \draw[book arrow] (a1)--(u);
  \draw[book arrow] (an)--(u);
\end{tikzpicture}
\caption{首次越界时刻把最大事件分成互不相交的 $\mathcal F_k$-可测区域；每条越界路径只被计算一次。}
\label{fig:6-2-2-first-crossing}
\end{figure}

弱型 $L^1$ 估计控制每一个尾概率；当 $p>1$ 时，层蛋糕公式可以把全部尾概率重新积成强
$L^p$ 估计。

\begin{theorem}[Doob $L^p$ 不等式]\label{theorem:6-2-2-2}
令 $p>1$，$q=p/(p-1)$。若 $M=(M_k)_{0\le k\le n}$ 是鞅且
$M_n\in L^p$，则
\[
  \|M_n^*\|_p\le q\|M_n\|_p.
\]
若 $M=(M_k)_{0\le k\le n}$ 是非负次鞅且 $M_n\in L^p$，同一结论成立，此时
$M_n^*=\max_{k\le n}M_k$。
\end{theorem}

\begin{proof}
在鞅情形，时间一致性与条件模长不等式给
\[
 |M_k|^p
 \le\bigl(\E[|M_n|\mid\mathcal F_k]\bigr)^p
 \le\E[|M_n|^p\mid\mathcal F_k],
\]
故 $M_k\in L^p$。非负次鞅情形由
$0\le M_k\le\E[M_n\mid\mathcal F_k]$ 和同一标量条件 Jensen 估计得到相同结论。
鞅情形令 $X_k=|M_k|$；条件模长不等式给
\[
 \E[X_{k+1}\mid\mathcal F_k]
 \ge \left|\E[M_{k+1}\mid\mathcal F_k]\right|=|M_k|,
\]
所以 $X$ 是非负次鞅。非负次鞅情形直接令 $X=M$。统一记
$X_n^*=\max_{k\le n}X_k$。对 $R>0$，尾积分公式与
定理~\ref{theorem:6-2-2-1} 给出
\begin{align*}
 \E[(X_n^*\wedge R)^p]
 &=p\int_0^R\lambda^{p-1}\Prob(X_n^*\ge\lambda)\dd\lambda\\
 &\le p\int_0^R\lambda^{p-2}
       \E[X_n\1_{\{X_n^*\ge\lambda\}}]\dd\lambda\\
 &=\frac{p}{p-1}\E\!\left[X_n(X_n^*\wedge R)^{p-1}\right].
\end{align*}
最后一式由 Tonelli 定理逐点积分得到。\Holder{} 不等式于是给出
\[
 \E[(X_n^*\wedge R)^p]
 \le q\|X_n\|_p
       \bigl\|(X_n^*\wedge R)^{p-1}\bigr\|_q
 =q\|X_n\|_p
       \|X_n^*\wedge R\|_p^{p-1}.
\]
若左边为零结论成立；否则约去其 $(p-1)/p$ 次幂，得到
$\|X_n^*\wedge R\|_p\le q\|X_n\|_p$。令 $R\uparrow\infty$，由单调收敛定理
得到结论。
\end{proof}

\begin{example}[随机游走最大值]\label{ex:6-2-2-1}
对方差为一的中心随机游走 $S_n=\sum_{k=1}^n\xi_k$，取 $p=2$ 得
\[
 \Prob\!\left(\max_{k\le n}|S_k|\ge\lambda\right)
 \le \frac{\E(S_n^*)^2}{\lambda^2}
 \le \frac{4\E S_n^2}{\lambda^2}
 =\frac{4n}{\lambda^2}.
\]
逐时刻使用 Chebyshev 不等式再作并集估计只给出
$\lambda^{-2}\sum_{k=1}^n\E S_k^2=n(n+1)/(2\lambda^2)$。最大不等式利用了路径的共同历史，
因此把二次量级降到一次量级。
\end{example}

接下来把振荡翻译成一项可预测交易。策略在 $X\le a$ 时买入一单位，随后一直持有，直到
$X\ge b$ 时卖出；完成一次交易便完成一次上穿。

\begin{theorem}[Doob 上穿不等式]\label{theorem:6-2-2-3}
若 $(X_k)_{0\le k\le n}$ 是实值次鞅，$a<b$，则
\begin{equation}\label{eq:6-2-upcrossing}
 (b-a)\E U_n[a,b]
 \le \E[(X_n-a)^-]+\E[X_n-X_0],
\end{equation}
其中 $x^- =\max\{-x,0\}$。
\end{theorem}

\begin{proof}
逐路径递归定义 $H_k\in\{0,1\}$。起初不持有；若在时刻 $k-1$ 不持有且
$X_{k-1}\le a$，则令 $H_k=1$ 并开始持有；持有期间令 $H_k=1$，直到某个时刻
$k-1$ 已满足 $X_{k-1}\ge b$，从下一增量起令 $H_k=0$；随后重复。这个规则只使用
$X_0,\ldots,X_{k-1}$，所以 $H_k$ 对 $\mathcal F_{k-1}$ 可测。

记 $G_n=\sum_{k=1}^nH_k(X_k-X_{k-1})$。每个已经卖出的单位至少赚 $b-a$。若时刻
$n$ 尚有一笔未完成持仓，买入价不高于 $a$，其账面收益至少为 $X_n-a$。因此逐路径有
\begin{equation}\label{eq:6-2-upcross-pathwise}
  (b-a)U_n[a,b]\le G_n+(X_n-a)^-.
\end{equation}
上述贪心规则完成的交易数确实等于定义
\ref{def:6-2-2-2} 中的最大上穿数。设它依次选出
\[
 s_1=\min\{k\le n:X_k\le a\},\qquad
 t_1=\min\{k>s_1:X_k\ge b\},
\]
并在 $t_j<n$ 时递归地取
\[
 s_{j+1}=\min\{k>t_j:X_k\le a\},\qquad
 t_{j+1}=\min\{k>s_{j+1}:X_k\ge b\},
\]
空集的最小值视为 $+\infty$。若这个规则完成了 $q$ 次，而
$s'_1<t'_1<\cdots<s'_m<t'_m$ 是任意一组上穿，则归纳得
\[
 s_j\le s'_j,\qquad t_j\le t'_j\qquad(1\le j\le m).
\]
其中 $j=1$ 来自最小值的定义；若对 $j$ 成立，则
$t_j\le t'_j<s'_{j+1}$，所以 $s_{j+1}\le s'_{j+1}$，进而
$t_{j+1}\le t'_{j+1}$。因此 $m\le q$；反过来，贪心规则自身给出长度为 $q$ 的上穿族，故
$q=U_n[a,b]$。于是前述收益分解正好给出
式~\eqref{eq:6-2-upcross-pathwise}。
由于 $H$ 非负且可预测，次鞅差分给 $\E G_n\ge0$，但这不是
\eqref{eq:6-2-upcrossing} 所需的方向。注意互补策略 $1-H$ 也非负且可预测，并且
\[
  X_n-X_0=G_n+
  \sum_{k=1}^n(1-H_k)(X_k-X_{k-1}).
\]
第二项期望非负，故 $\E G_n\le\E(X_n-X_0)$。对
\eqref{eq:6-2-upcross-pathwise} 取期望即得结论。
\end{proof}

末项 $\E[X_n-X_0]$ 不能随意换成只含 $X_0$ 正部的表达式。例如取确定性递增次鞅
$X_k=k$ 及 $0<a<b<n$，上穿次数为一；若丢掉终端漂移，右端可能变成零。公式的方向与
定义~\ref{def:6-2-2-2} 所数的“低到高”上穿严格配套。

\begin{theorem}[鞅／次鞅收敛定理]\label{theorem:6-2-2-4}
若实值次鞅 $(X_n)$ 满足
\[
  \sup_{n\ge0}\E X_n^+<\infty,
\]
则存在有限的 $X_\infty\in L^1$，使 $X_n\to X_\infty$ a.s.；这里不声称
$L^1$ 收敛。

若鞅 $(M_n)$ 对某个 $p>1$ 满足 $\sup_n\|M_n\|_p<\infty$，则存在
$M_\infty\in L^p$，使 $M_n\to M_\infty$ a.s. 且在 $L^p$ 中收敛。
\end{theorem}

\begin{proof}
令 $C=\sup_n\E X_n^+<\infty$。对固定有理数 $a<b$，
定理~\ref{theorem:6-2-2-3} 的右端满足
\begin{align*}
 \E[(X_n-a)^-]+\E[X_n-X_0]
 &=\E[\max\{X_n,a\}]-\E X_0\\
 &\le C+a^+-\E X_0.
\end{align*}
因此 $\sup_n\E U_n[a,b]<\infty$。由于 $U_n[a,b]\uparrow U_\infty[a,b]$，
单调收敛定理给 $\E U_\infty[a,b]<\infty$，从而
$U_\infty[a,b]<\infty$ a.s. 对每一对有理数 $a<b$ 成立。这样的有理数对只有可数多个，
故可取一个共同的概率一事件。

若在该事件上 $\liminf_nX_n<\limsup_nX_n$，可在两者之间选有理数 $a<b$；数列随后会完成
无穷多次从 $a$ 以下到 $b$ 以上的上穿，与 $U_\infty[a,b]<\infty$ 矛盾。因此
$X_n$ 在扩展实数中有极限。又因次鞅的均值非降，
\[
 \E X_n^- =\E X_n^+-\E X_n\le C-\E X_0,
\]
所以 $\sup_n\E|X_n|\le2C-\E X_0<\infty$。Fatou 引理给
\[
 \E\!\left[\liminf_n|X_n|\right]
 \le\liminf_n\E|X_n|<\infty.
\]
这排除了极限为 $\pm\infty$ 的正概率事件，并且给出
$X_\infty\in L^1$。

现在设 $M$ 在 $L^p$ 中有界。若 $M$ 为实值，它也在 $L^1$ 中有界，故第一部分给出 a.s. 极限；
若 $M$ 为复值，则 $\Re M$ 与 $\Im M$ 都是 $L^1$ 有界的实鞅，分别应用第一部分后仍得到
$M_n\to M_\infty$ a.s.。对每个 $N$，定理~\ref{theorem:6-2-2-2} 给
\[
 \E\!\left[\max_{k\le N}|M_k|^p\right]
 \le q^p\sup_n\E|M_n|^p,
 \qquad q=\frac{p}{p-1}.
\]
令 $N\to\infty$，得到 $M^*=\sup_k|M_k|\in L^p$。于是
$|M_n-M_\infty|^p\le(2M^*)^p$，且左边 a.s. 趋于零；支配收敛定理给
$\|M_n-M_\infty\|_p\to0$，并同时说明 $M_\infty\in L^p$。
\end{proof}

上一定理的第一部分没有给 $L^1$ 收敛；缺失的条件正是一致可积。

\begin{proposition}[UI、$L^1$ 收敛与闭合]\label{proposition:6-2-2-5}
对离散时间鞅 $(M_n)$，下列三项等价：
\begin{enumerate}
 \item 随机变量族 $\{M_n:n\ge0\}$ 一致可积；
 \item 存在 $M_\infty\in L^1$，使 $M_n\to M_\infty$ 于 $L^1$；
 \item 存在 $Y\in L^1$，使 $M_n=\E[Y\mid\mathcal F_n]$ 对所有 $n$ 成立。
\end{enumerate}
在第 (2) 项成立时，可在第 (3) 项取 $Y=M_\infty$。
\end{proposition}

\begin{proof}
设 (1) 成立。一致可积蕴含 $\sup_n\E|M_n|<\infty$。实值情形由
定理~\ref{theorem:6-2-2-4} 得到 a.s. 收敛；复值情形分别把该定理用于实部和虚部。
因此 $M_n$ a.s. 收敛到某个 $M_\infty\in L^1$。
Vitali 定理~\ref{theorem:5-3-2-1} 把 a.s. 收敛与 UI 合并，得到
$M_n\to M_\infty$ 于 $L^1$，所以 (2) 成立。

设 (2) 成立。固定 $n$。对每个 $m\ge n$，时间一致性给
$M_n=\E[M_m\mid\mathcal F_n]$。条件期望的 $L^1$ 收缩性给
\[
 \bigl\|\E[M_m-M_\infty\mid\mathcal F_n]\bigr\|_1
 \le\|M_m-M_\infty\|_1\longrightarrow0.
\]
因此 $M_n=\E[M_\infty\mid\mathcal F_n]$ a.s.，得到 (3)。

最后设 (3) 成立，并令 $A_{n,K}=\{|M_n|>K\}\in\mathcal F_n$。由条件模长
不等式，
\[
 \E[|M_n|\1_{A_{n,K}}]
 \le\E[|Y|\1_{A_{n,K}}].
\]
而 $\|M_n\|_1\le\|Y\|_1$，故
$\Prob(A_{n,K})\le\|Y\|_1/K$。对任意 $R>0$，
\[
 \E[|Y|\1_{A_{n,K}}]
 \le \E[|Y|\1_{\{|Y|>R\}}]
      +R\frac{\|Y\|_1}{K}.
\]
先取 $R$ 使第一项小，再令 $K\to\infty$，所得界与 $n$ 无关，故
$\{M_n\}$ UI。
\end{proof}

\begin{example}[非 UI 非负鞅]\label{ex:6-2-2-2}
考虑 Galton--Watson 过程：$Z_0=1$，每个个体独立地产生 $0$ 或 $2$ 个后代，概率各为
$1/2$。给定第 $n$ 代的信息 $\mathcal F_n$，每个现存个体的平均后代数为一，所以
\[
  \E[Z_{n+1}\mid\mathcal F_n]=Z_n,
  \qquad \E Z_n=1.
\]
令 $q_n=\Prob(Z_n=0)$。从一个祖先出发，下一代要么立即灭绝，要么由两个独立子过程组成，故
\[
  q_0=0,\qquad q_{n+1}=\frac{1+q_n^2}{2}.
\]
由 $0\le q_n\le1$ 及
\[
 q_{n+1}-q_n=\frac{(1-q_n)^2}{2}\ge0,
\]
归纳得序列 $(q_n)$ 递增且不超过一；其极限 $q$ 满足
$q=(1+q^2)/2$，即 $(q-1)^2=0$，所以 $q_n\uparrow1$。灭绝后过程保持为零，因而
\[
 \Prob(\text{最终灭绝})=\Prob\Bigl(\bigcup_n\{Z_n=0\}\Bigr)=\lim_nq_n=1,
\]
从而 $Z_n\to0$ a.s.。但 $\E Z_n=1$ 对所有 $n$ 成立，所以不可能在 $L^1$ 收敛到零；
命题~\ref{proposition:6-2-2-5} 说明 $(Z_n)$ 不是 UI。稀有而庞大的幸存族群携带了没有出现在
点态极限中的全部期望。
\end{example}

\begin{example}[条件期望逼近]\label{ex:6-2-2-3}
设 $X\in L^p$、$p>1$，并令 $M_n=\E[X\mid\mathcal F_n]$。条件 Jensen 给
$\|M_n\|_p\le\|X\|_p$，故定理~\ref{theorem:6-2-2-4} 给出 a.s. 及 $L^p$ 极限
$M_\infty$。若
\[
  \mathcal F_\infty=\sigma\Bigl(\bigcup_{n\ge0}\mathcal F_n\Bigr),
\]
则 $M_\infty$ 对 $\mathcal F_\infty$ 可测。对 $A\in\bigcup_n\mathcal F_n$，从某一时刻起
$\E[M_k\1_A]=\E[X\1_A]$；令 $k\to\infty$ 得同一等式对 $M_\infty$ 成立。
包含这些 $A$ 且保持该等式的事件类是 Dynkin 系统，而
$\bigcup_n\mathcal F_n$ 是代数，故 $\pi$--$\lambda$ 定理推出
\[
  M_\infty=\E[X\mid\mathcal F_\infty].
\]
若 $X$ 本身对 $\mathcal F_\infty$ 可测，极限才等于 $X$；“信息递增”不等于“最终揭示一切”。
\end{example}

\begin{figure}[htbp]
\centering
\begin{tikzpicture}[x=0.72cm,y=0.72cm]
  \draw[->,BookInk!75] (0,0)--(12.2,0) node[right] {$k$};
  \draw[->,BookInk!75] (0,-1.2)--(0,5.2) node[above] {$X_k$};
  \draw[BookAccent,semithick] (0,1.2)--(12,1.2) node[right] {$a$};
  \draw[BookAccent,semithick] (0,3.6)--(12,3.6) node[right] {$b$};
  \draw[BookInk,thick]
    plot coordinates {(0,2.4) (1,0.8) (2,1.0) (3,2.0) (4,3.9)
      (5,3.0) (6,0.7) (7,1.1) (8,2.5) (9,4.1) (10,2.7) (11,0.9) (12,2.0)};
  \fill[BookWarm] (1,0.8) circle (2pt) node[below left] {$s_1$};
  \fill[BookWarm] (4,3.9) circle (2pt) node[above] {$t_1$};
  \fill[BookWarm] (6,0.7) circle (2pt) node[below] {$s_2$};
  \fill[BookWarm] (9,4.1) circle (2pt) node[above] {$t_2$};
  \draw[decorate,decoration={brace,amplitude=4pt},BookWarm]
    (1.1,0.25)--(3.9,0.25) node[midway,below=5pt] {持有};
  \draw[decorate,decoration={brace,amplitude=4pt},BookWarm]
    (6.1,0.25)--(8.9,0.25) node[midway,below=5pt] {持有};
\end{tikzpicture}
\caption{可预测策略只在看到 $X\le a$ 后买入，在首次达到 $X\ge b$ 后卖出；两笔完成交易对应两次上穿。}
\label{fig:6-2-2-upcrossings}
\end{figure}

上穿次数有限排除的是实序列的无限振荡；极限为 $\pm\infty$ 仍需用 $L^1$ 界排除，这正是
定理~\ref{theorem:6-2-2-4} 证明的最后一步。另一方面，$L^1$ 有界鞅未必在 $L^1$ 中收敛，
例~\ref{ex:6-2-2-2} 已给出完全可计算的反例。$p=1$ 处最大不等式只剩
\eqref{eq:6-2-doob-weak-one} 的弱型估计；强型常数 $p/(p-1)$ 在 $p\downarrow1$ 时发散并非偶然。

\begin{exercise}
从定理~\ref{theorem:6-2-2-1} 出发，逐行重做定理~\ref{theorem:6-2-2-2} 的尾积分证明，
并解释为何 $p=1$ 时出现不可积的因子 $\lambda^{-1}$。
\end{exercise}

\begin{exercise}
对一个确定性实数列写出生成 $H_k\in\{0,1\}$ 的递归规则，证明
\eqref{eq:6-2-upcross-pathwise}，并验证 $H_k$ 只依赖 $x_0,\ldots,x_{k-1}$。
\end{exercise}

\begin{exercise}
证明若 $\sup_n\|M_n\|_p<\infty$、$p>1$，则 $\{M_n:n\ge0\}$ 一致可积。
比较这一短证与定理~\ref{theorem:6-2-2-4} 中得到 $L^p$ 收敛的证明。
\end{exercise}

有了路径最大值与极限控制，下一步可以在由路径自身决定的时刻读取鞅。但“看见赢了就停”只有在
停下的决定不借用未来，并且尾部质量不会逃到无穷时才安全。

\subsection{停时、可选抽样与一致可积性}\label{subsec:6-2-3}
\index{停时}\index{可选抽样}\index{停止过程}

“第一次达到十元”可由目前为止的记录判断；“最后一次达到十元”通常要等未来全部结束以后才知道。
这个差别决定了随机时刻能否与条件期望相容。

\begin{definition}[停时]\label{def:6-2-3-1}
取值于 $\{0,1,2,\ldots\}\cup\{\infty\}$ 的随机变量 $\tau$ 若满足
\[
  \{\tau\le n\}\in\mathcal F_n\qquad(n\ge0),
\]
则称 $\tau$ 为关于滤过 $(\mathcal F_n)$ 的停时。等价地，
$\{\tau=n\}\in\mathcal F_n$ 对每个有限 $n$ 成立。
\end{definition}

等价性来自
\[
 \{\tau=n\}=\{\tau\le n\}\setminus\{\tau\le n-1\},
 \qquad
 \{\tau\le n\}=\bigcup_{k=0}^n\{\tau=k\},
\]
其中 $n=0$ 时把第二个被减集合视为空集。特别地，
$\{\tau\ge k\}=\{\tau\le k-1\}^{\mathrm c}\in\mathcal F_{k-1}$：在第 $k$ 个增量出现前，
我们知道过程是否还没有停。

\begin{definition}[停止过程与停止信息]\label{def:6-2-3-2}
过程 $X$ 在停时 $\tau$ 处的停止过程为
\[
  X_n^\tau=X_{n\wedge\tau}.
\]
停下时可知的事件构成
\[
 \mathcal F_\tau
 =\{A\in\mathcal F:A\cap\{\tau\le n\}\in\mathcal F_n
                  \text{ 对每个 }n\}.
\]
这通常远大于 $\sigma(\tau)$：它不仅知道何时停，也保留停下以前观察到的路径信息。
\end{definition}

\begin{lemma}[停止信息的基本性质]\label{lemma:6-2-3-stopped-sigma}
$\mathcal F_\tau$ 是 $\sigma$-代数。若 $\sigma\le\tau$ 是两个停时，则
$\mathcal F_\sigma\subset\mathcal F_\tau$。若 $X$ 适应且 $\tau<\infty$，则
$X_\tau$ 对 $\mathcal F_\tau$ 可测；若 $\tau\le N$，则 $X_\tau$ 可积，只要
$X_0,\ldots,X_N$ 均可积。
\end{lemma}

\begin{proof}
空集属于 $\mathcal F_\tau$。若 $A\in\mathcal F_\tau$，则
\[
 A^{\mathrm c}\cap\{\tau\le n\}
 =\{\tau\le n\}\setminus(A\cap\{\tau\le n\})\in\mathcal F_n.
\]
可数并的条件由交与并的分配律逐项得到，故 $\mathcal F_\tau$ 是 $\sigma$-代数。

若 $A\in\mathcal F_\sigma$ 且 $\sigma\le\tau$，则
\[
 A\cap\{\tau\le n\}
 =(A\cap\{\sigma\le n\})\cap\{\tau\le n\}\in\mathcal F_n,
\]
所以 $A\in\mathcal F_\tau$。最后，对 Borel 集 $B$，
\[
 \{X_\tau\in B\}\cap\{\tau\le n\}
 =\bigcup_{k=0}^n\bigl(\{\tau=k\}\cap\{X_k\in B\}\bigr)\in\mathcal F_n,
\]
这证明停止值可测。若 $\tau\le N$，则
$|X_\tau|\le\sum_{k=0}^N|X_k|$，故可积。
\end{proof}

\begin{example}[首次命中]\label{ex:6-2-3-1}
若 $X$ 适应，$A$ 是状态空间中的可测集，则
\[
  \tau_A=\inf\{n\ge0:X_n\in A\}
\]
是停时，因为
\[
  \{\tau_A\le n\}=\bigcup_{k=0}^n\{X_k\in A\}\in\mathcal F_n.
\]
约定空集的下确界为 $\infty$，所以从未命中也已包含在定义中。
\end{example}

\begin{example}[最后命中不是停时]\label{ex:6-2-3-2}
令 $S_0=0$，$S_n$ 为简单对称随机游走，并只观察到时刻 $2$。定义
\[
  L=\max\{k\in\{0,1,2\}:S_k=0\}.
\]
$S_1$ 必为 $\pm1$；若第二步与第一步反向，则 $L=2$，否则 $L=0$。因此
$\Prob(L\le0)=1/2$。自然滤过的 $\mathcal F_0$ 是平凡 $\sigma$-代数，而
$\{L\le0\}$ 概率既非零也非一，所以该事件不在 $\mathcal F_0$，$L$ 不是停时。
失败的不是“最大值”这个运算，而是判定最后一次命中必须查看未来。
\end{example}

\begin{definition}[可选抽样问题]\label{def:6-2-3-3}
对停时 $\sigma\le\tau$，可选抽样问题询问何时有
\[
  \E[M_\tau\mid\mathcal F_\sigma]=M_\sigma.
\]
该式首先要求停止值可测且可积，还要求从有界截断
$\sigma\wedge n,\tau\wedge n$ 放开时能够交换极限与期望。
\end{definition}

\begin{proposition}[停止保持鞅性]\label{proposition:6-2-3-1}
若 $M$ 是鞅、$\tau$ 是停时，则停止过程 $(M_{n\wedge\tau})_{n\ge0}$ 仍是鞅。
\end{proposition}

\begin{proof}
由 $\{\tau\ge k\}\in\mathcal F_{k-1}$，过程
$H_k=\1_{\{\tau\ge k\}}$ 有界且可预测。逐路径有
\[
 M_{n\wedge\tau}
 =M_0+\sum_{k=1}^n\1_{\{\tau\ge k\}}(M_k-M_{k-1}).
\]
命题~\ref{proposition:6-2-1-4} 因而说明右边是鞅。
\end{proof}

\begin{theorem}[有界可选抽样]\label{theorem:6-2-3-2}
若 $\sigma\le\tau\le N$ 是有界停时，$M$ 是鞅，则
\[
  \E[M_\tau\mid\mathcal F_\sigma]=M_\sigma\quad\text{a.s.}
\]
特别地，$\E M_\tau=\E M_\sigma=\E M_0$。
\end{theorem}

\begin{proof}
停止值可积且 $M_\sigma$ 对 $\mathcal F_\sigma$ 可测。固定
$A\in\mathcal F_\sigma$。有限和恒等式
\[
 M_\tau-M_\sigma
 =\sum_{k=1}^N\1_{\{\sigma<k\le\tau\}}(M_k-M_{k-1})
\]
逐路径成立。对每个 $k$，
\[
 A\cap\{\sigma<k\le\tau\}
 =(A\cap\{\sigma\le k-1\})\cap\{\tau\ge k\}
 \in\mathcal F_{k-1}.
\]
因此
\begin{align*}
 \E[\1_A\1_{\{\sigma<k\le\tau\}}(M_k-M_{k-1})]
 &=\E\!\left[\1_A\1_{\{\sigma<k\le\tau\}}
     \E[M_k-M_{k-1}\mid\mathcal F_{k-1}]\right]\\
 &=0.
\end{align*}
对有限个 $k$ 求和，得到 $\E[M_\tau\1_A]=\E[M_\sigma\1_A]$。这正是所列条件
期望等式的积分刻画。取 $A=\Omega$ 得 $\E M_\tau=\E M_\sigma$；再把已经证明的条件式
用于停时对 $0\le\sigma$，得到 $\E M_\sigma=\E M_0$。
\end{proof}

有界性允许逐项处理有限和。对无界停时，正确的替代条件不是“几乎处处有限”本身，而是保证全部
截断停止值的尾部不会逃逸。

\begin{theorem}[UI 可选抽样]\label{theorem:6-2-3-3}
设 $M$ 是一致可积鞅，$\sigma\le\tau<\infty$ a.s. 是停时。则
$M_\sigma,M_\tau\in L^1$，并且
\[
  \E[M_\tau\mid\mathcal F_\sigma]=M_\sigma\quad\text{a.s.}
\]
\end{theorem}

\begin{proof}
由命题~\ref{proposition:6-2-2-5}，存在 $M_\infty\in L^1$ 使
\[
  M_n=\E[M_\infty\mid\mathcal F_n],
  \qquad M_n\longrightarrow M_\infty\quad\text{于 }L^1\text{ 及 a.s.}
\]
先取任意有界停时 $\rho\le N$。对 $A\in\mathcal F_\rho$，
定理~\ref{theorem:6-2-3-2} 与
$M_N=\E[M_\infty\mid\mathcal F_N]$ 给出
\[
 \E[M_\rho\1_A]=\E[M_N\1_A]=\E[M_\infty\1_A].
\]
故 $M_\rho=\E[M_\infty\mid\mathcal F_\rho]$。令
$B_{\rho,K}=\{|M_\rho|>K\}$。条件模长不等式、Markov 不等式和任意 $R>0$ 给出
\[
 \E[|M_\rho|\1_{B_{\rho,K}}]
 \le \E[|M_\infty|\1_{\{|M_\infty|>R\}}]
      +R\frac{\|M_\infty\|_1}{K}.
\]
先取 $R$ 使第一项任意小，再令 $K\to\infty$；该界与 $\rho$ 无关。所以
\[
 \{M_\rho:\rho\text{ 为有界停时}\}
\]
是一致可积族。

令 $\rho_n=\tau\wedge n$。因为 $\tau<\infty$ a.s.，
$M_{\rho_n}\to M_\tau$ a.s.；上面的 UI 与 Vitali 定理给
$M_{\rho_n}\to M_\tau$ 于 $L^1$，特别 $M_\tau$ 可积。若
$A\in\mathcal F_\tau$，置 $A_n=A\cap\{\tau\le n\}$。若 $m<n$，则
\[
 A_n\cap\{\rho_n\le m\}=A\cap\{\tau\le m\}\in\mathcal F_m;
\]
若 $m\ge n$，则该交集为 $A_n\in\mathcal F_n\subset\mathcal F_m$。故
$A_n\in\mathcal F_{\rho_n}$，于是
\[
 \E[M_{\rho_n}\1_{A_n}]=\E[M_\infty\1_{A_n}].
\]
令 $n\to\infty$：左边用 $L^1$ 收敛，右边用 DCT，得到
$\E[M_\tau\1_A]=\E[M_\infty\1_A]$，即
\[
 M_\tau=\E[M_\infty\mid\mathcal F_\tau].
\]
令 $\eta_n=\sigma\wedge n$。有界停时公式给
$M_{\eta_n}=\E[M_\infty\mid\mathcal F_{\eta_n}]$。因为 $\sigma<\infty$ a.s.，
$M_{\eta_n}\to M_\sigma$ a.s.；有界停时停止值族的统一可积性与 Vitali 定理再给
$L^1$ 收敛。固定 $B\in\mathcal F_\sigma$，置 $B_n=B\cap\{\sigma\le n\}$。与 $A_n$ 的核验
相同，$B_n\in\mathcal F_{\eta_n}$，从而
\[
 \E[M_{\eta_n}\1_{B_n}]=\E[M_\infty\1_{B_n}].
\]
当 $n\to\infty$ 时，左边由 $L^1$ 收敛及
$\E[|M_\sigma|\1_{B\setminus B_n}]\to0$ 收敛到
$\E[M_\sigma\1_B]$；右边由 DCT 收敛到 $\E[M_\infty\1_B]$。故
$M_\sigma=\E[M_\infty\mid\mathcal F_\sigma]$。由
引理~\ref{lemma:6-2-3-stopped-sigma}，
$\mathcal F_\sigma\subset\mathcal F_\tau$；塔式性质因此给出所求等式。
\end{proof}

若只允许停时在一个零集上取 $\infty$，可在该零集上定义
$M_\tau=M_\infty$，上述结论不变。若 $\Prob(\tau=\infty)>0$，则还必须规定“无穷时刻的信息”。
一种安全约定是令
$\mathcal F_\infty=\sigma(\bigcup_n\mathcal F_n)$，并在
定义~\ref{def:6-2-3-2} 的条件之外要求
$A\cap\{\tau=\infty\}\in\mathcal F_\infty$，同时在该事件上置
$M_\tau=M_\infty$。若不加这项终端条件，原定义对 $\{\tau=\infty\}$ 上的集合没有任何约束，
不能无条件声称 $M_\tau$ 对停止信息可测。这是“允许取 $\infty$”时不能省略的边界。

\begin{corollary}[可积停时与有界增量]\label{corollary:6-2-3-bounded-increments}
设 $M$ 是鞅，并且存在常数 $C<\infty$ 使
$|M_k-M_{k-1}|\le C$ a.s. 对所有 $k$ 成立。若停时
$\sigma\le\tau$ 满足 $\E\tau<\infty$，则 $M_\sigma,M_\tau\in L^1$，且
\[
  \E[M_\tau\mid\mathcal F_\sigma]=M_\sigma.
\]
\end{corollary}

\begin{proof}
由 $\sigma\le\tau$ 与 $\E\tau<\infty$，两个停时都 a.s. 有限，而且
\[
 |M_\sigma|\le|M_0|+C\sigma\le|M_0|+C\tau,
 \qquad |M_\tau|\le|M_0|+C\tau,
\]
故两个停止值可积。
由命题~\ref{proposition:6-2-3-1}，$N_n=M_{n\wedge\tau}$ 是鞅。又因
\[
 |N_n-M_\tau|
 \le C(\tau-n)^+,
\]
而 $(\tau-n)^+\downarrow0$ 且被可积变量 $\tau$ 支配，DCT 给
$N_n\to M_\tau$ 于 $L^1$。命题~\ref{proposition:6-2-2-5} 因而说明 $N$ 是 UI 鞅。
由于 $\sigma\le\tau$，有 $N_\sigma=M_\sigma$、$N_\tau=M_\tau$；对 $N$ 应用
定理~\ref{theorem:6-2-3-3} 即得结论。
\end{proof}

\begin{proposition}[随机游走赌徒破产]\label{proposition:6-2-3-4}
令 $(S_n)$ 为从 $i\in\{0,\ldots,N\}$ 出发的简单对称随机游走，并令
\[
  \tau=\inf\{n\ge0:S_n\in\{0,N\}\}.
\]
则
\[
  \Prob_i(S_\tau=N)=\frac{i}{N},
  \qquad
  \E_i\tau=i(N-i).
\]
\end{proposition}

\begin{proof}
边界起点的结论直接成立，以下设 $0<i<N$。令 $\tau_m=\tau\wedge m$。过程
$S_n$ 与 $S_n^2-n$ 都是鞅：后一结论是命题~\ref{proposition:6-2-1-5} 后平方随机游走计算中
$\sigma^2=1$ 的情形。对有界停时 $\tau_m$ 使用
定理~\ref{theorem:6-2-3-2}，得到
\begin{equation}\label{eq:6-2-gambler-truncation}
  \E_i S_{\tau_m}=i,
  \qquad
  \E_i S_{\tau_m}^2-\E_i\tau_m=i^2.
\end{equation}
在命中以前 $S_n\in\{1,\ldots,N-1\}$，命中时位于 $0$ 或 $N$，故
$0\le S_{\tau_m}\le N$。由第二式，
\[
  0\le\E_i\tau_m=\E_iS_{\tau_m}^2-i^2\le N^2.
\]
单调收敛定理先给 $\E_i\tau\le N^2<\infty$，因而 $\tau<\infty$ a.s.。现在才放开截断。
由 $S_{\tau_m}\to S_\tau$ a.s. 及统一界 $N$，DCT 和
\eqref{eq:6-2-gambler-truncation} 给 $\E_iS_\tau=i$。若
$p=\Prob_i(S_\tau=N)$，则 $S_\tau$ 只取 $0,N$，所以 $Np=i$，即 $p=i/N$。

最后，对第二式令 $m\to\infty$；左边的停时项用 MCT，平方项用 DCT，得到
\[
 \E_i\tau=\E_iS_\tau^2-i^2=N^2\frac{i}{N}-i^2=i(N-i).
\]
证明次序很重要：$\E\tau<\infty$ 是先由有界截断推出的，不是预先借用待证公式。
\end{proof}

\begin{proposition}[可积停时下的 Wald 第一恒等式]\label{proposition:6-2-3-5}
设 $\xi_1,\xi_2,\ldots$ i.i.d.，$\E|\xi_1|<\infty$，
$\E\xi_1=\mu$，并令
$\mathcal F_n=\sigma(\xi_1,\ldots,\xi_n)$、
$S_n=\sum_{k=1}^n\xi_k$。若停时 $\tau$ 满足 $\E\tau<\infty$，则
$S_\tau=\sum_{k=1}^{\tau}\xi_k$ 可积，且
\[
  \E S_\tau=\mu\E\tau.
\]
\end{proposition}

\begin{proof}
事件 $\{\tau\ge k\}$ 属于 $\mathcal F_{k-1}$，而 $\xi_k$ 独立于
$\mathcal F_{k-1}$。Tonelli 定理给出
\begin{align*}
 \E\sum_{k=1}^{\tau}|\xi_k|
 &=\sum_{k=1}^\infty\E[\1_{\{\tau\ge k\}}|\xi_k|]\\
 &=\E|\xi_1|\sum_{k=1}^\infty\Prob(\tau\ge k)
 =\E|\xi_1|\E\tau<\infty.
\end{align*}
所以随机级数绝对可积，可以交换求和与期望。再次利用独立性，
\begin{align*}
 \E S_\tau
 &=\sum_{k=1}^\infty
   \E[\1_{\{\tau\ge k\}}\xi_k]
 =\sum_{k=1}^\infty \Prob(\tau\ge k)\E\xi_k
 =\mu\E\tau.
\end{align*}
\end{proof}

\begin{example}[第 $r$ 次 Bernoulli 成功]
令 $\xi_k\sim\mathrm{Bernoulli}(p)$，$0<p<1$，取整数 $r\ge1$，并令
$\tau=\inf\{n:S_n=r\}$。使用 Wald 恒等式以前先验证 $\E\tau<\infty$。
把试验按长度 $r$ 分块；每一整块全部成功的概率为 $p^r$，各块相互独立。若
$\tau>mr$，则前 $m$ 块中没有一块全部成功，故
\[
  \Prob(\tau>mr)\le(1-p^r)^m.
\]
把尾和按这些长度为 $r$ 的块分组，得到
\[
 \E\tau=\sum_{k\ge1}\Prob(\tau\ge k)
 \le r\sum_{m\ge0}(1-p^r)^m=\frac{r}{p^r}<\infty.
\]
命中时 $S_\tau=r$，命题~\ref{proposition:6-2-3-5} 因而给
$r=p\E\tau$，即 $\E\tau=r/p$。粗尾界只负责合法化 Wald；最后的精确答案来自恒等式。
当 $p=1$ 时 $\tau=r$，同一答案可直接读出。
\end{example}

最后看一个专门揭示假设缺口的策略。

\begin{example}[加倍下注策略的失效]\label{ex:6-2-3-3}
令公平硬币结果 $\varepsilon_k$ 在赢时为 $1$、输时为 $-1$，并令
$\tau=\inf\{k\ge1:\varepsilon_k=1\}$。在第 $k$ 局尚未赢时下注
$2^{k-1}$，即取
\[
 H_k=2^{k-1}\1_{\{\tau\ge k\}},
 \qquad W_n=\sum_{k=1}^nH_k\varepsilon_k.
\]
$H_k$ 在第 $k$ 次抛掷前可知；对每个固定 $n$，差分条件均值为零，所以 $(W_n)$ 是鞅。
又因 $\Prob(\tau>n)=2^{-n}$，$\tau<\infty$ a.s.。若首次在第 $k$ 局赢，过去总损失为
$1+2+\cdots+2^{k-2}=2^{k-1}-1$，本局赢得 $2^{k-1}$，故
$W_\tau=1$ a.s.

这不与可选抽样矛盾。事实上
\[
 W_n=
 \begin{cases}
 1,&\tau\le n,\\
 -(2^n-1),&\tau>n,
 \end{cases}
 \qquad \E W_n=0.
\]
在概率 $2^{-n}$ 的尾事件上，损失大小接近 $2^n$；因此对任意 $K$，取足够大的 $n$ 有
\[
 \E[|W_n|\1_{\{|W_n|>K\}}]
 \ge (2^n-1)2^{-n}\longrightarrow1.
\]
所以 $(W_n)$ 不一致可积。a.s. 有限的停时并未阻止期望质量沿越来越稀有的失败路径逃到无穷。
\end{example}

\begin{figure}[htbp]
\centering
\begin{tikzpicture}[x=1.15cm,y=0.72cm]
  \draw[->,BookInk!75] (0,0)--(8.4,0) node[right] {时间};
  \foreach \x in {0,...,7}{
    \draw[BookInk!45] (\x,-0.12)--(\x,0.12);
    \node[below] at (\x,-0.12) {\scriptsize $\x$};
  }
  \fill[BookAccent!12] (2.05,0.35) rectangle (6.95,1.1);
  \node[book label] at (4.5,0.72) {$\{\sigma<k\le\tau\}$：第 $k$ 步前即可判定是否仍在抽样};
  \draw[BookAccent,thick] (2,0.25)--(2,1.25) node[above] {$\sigma$};
  \draw[BookAccent,thick] (7,0.25)--(7,1.25) node[above] {$\tau$};
  \begin{scope}[yshift=-1.65cm]
    \draw[BookWarm,thick] (0,0)--(8,0) node[right] {$0$};
    \draw[BookWarm,thick] (0,2)--(8,2) node[right] {$N$};
    \draw[BookInk,thick] plot coordinates
      {(0,0.8) (1,1.2) (2,0.8) (3,1.2) (4,1.6) (5,1.2) (6,1.6) (7,2)};
    \fill[BookWarm] (7,2) circle (2pt) node[above left] {首次命中};
  \end{scope}
\end{tikzpicture}
\caption{上图的随机区间给出可选抽样证明中的可预测系数；下图是停于 $0,N$ 的赌徒破产路径。}
\label{fig:6-2-3-stopping}
\end{figure}

有界停时、UI 鞅和 Wald 定理中的“可积停时加 i.i.d. 可积增量”分别是三组不同的充分条件。
它们虽然都导出停止值的期望恒等式，却不可相互替代。尤其，$\mathcal F_\tau$ 不是
$\sigma(\tau)$，停止值可积也不能由把 $n$ 换成 $\tau$ 的符号操作得到；每次放开截断都必须明确
使用 DCT、MCT 或 UI/Vitali 中的一种。

\begin{exercise}
证明命题~\ref{proposition:6-2-3-1} 时，不引用可预测变换，直接分别计算
$\{\tau\le n\}$ 与 $\{\tau>n\}$ 上的
$\E[M_{(n+1)\wedge\tau}\mid\mathcal F_n]$。
\end{exercise}

\begin{exercise}
设 $\sigma\le\tau\le N$。证明事件
$\{\sigma<k\le\tau\}$ 属于 $\mathcal F_{k-1}$，并用这一事实独立重证有界可选抽样。
\end{exercise}

\begin{exercise}
逐项复核命题~\ref{proposition:6-2-3-4}：先由平方鞅得到
$\sup_m\E\tau_m<\infty$，再说明两个极限分别使用 MCT 与 DCT。若随机游走每步向右概率为
$p\ne1/2$，指出 $S_n$ 的哪一个指数函数可能替代线性鞅。
\end{exercise}

\begin{exercise}
在例~\ref{ex:6-2-3-3} 中直接计算 $\E|W_n-W_\tau|$，说明 a.s. 收敛为什么不升级为
$L^1$ 收敛，并指出定理~\ref{theorem:6-2-3-3} 的哪项假设失败。
\end{exercise}

鞅把完整过去压缩成一个条件平均。下一节将研究另一种压缩：对 Markov 过程，给定当前状态以后，
完整过去不再改变未来的条件分布。转移核和半群正是把这种状态压缩组织成可迭代算子的语言。
